% !TEX program = xelatex
\documentclass[11pt,a4paper]{ctexart}

\usepackage{geometry}
\geometry{margin=1.7cm}

\usepackage{amsmath,amssymb,amsthm,mathtools,bm}
\usepackage{enumitem}
\usepackage{xcolor}
\usepackage{hyperref}
\usepackage{multicol}
\usepackage[most]{tcolorbox}
\usepackage{tikz}
\usepackage{microtype}
\usepackage{titlesec}
\usepackage{cleveref}
\usepackage{fvextra}

\setlength{\columnsep}{0.8cm}
\setlength{\parindent}{2em}
\setlength{\parskip}{0.3em}
\raggedcolumns

\definecolor{accent}{RGB}{39,113,172}
\definecolor{softback}{RGB}{246,249,253}
\definecolor{softgray}{RGB}{110,120,135}

\hypersetup{
  colorlinks=true,
  linkcolor=accent!80!black,
  urlcolor=accent!80!black,
  citecolor=accent!80!black
}

\setlist[itemize]{leftmargin=1.2em,itemsep=0.2em,topsep=0.25em}
\setlist[enumerate]{leftmargin=1.4em,itemsep=0.2em,topsep=0.25em}

\titleformat{\section}{\large\bfseries\color{accent}}{\thesection}{0.6em}{}[\titlerule]
\titleformat{\subsection}{\normalsize\bfseries\color{accent!90}}{\thesubsection}{0.5em}{}
\titleformat{\subsubsection}{\normalsize\bfseries\color{accent!80}}{\thesubsubsection}{0.5em}{}
\titlespacing*{\section}{0pt}{0.8em}{0.45em}
\titlespacing*{\subsection}{0pt}{0.6em}{0.3em}
\titlespacing*{\subsubsection}{0pt}{0.5em}{0.2em}

\tcbset{
  boxrule=0.55pt,
  arc=2pt,
  left=6pt,right=6pt,top=5pt,bottom=5pt,
  colframe=accent!25,
  colback=softback,
}

\newtcolorbox{examplebox}[1]{
  enhanced,
  colback=accent!3,
  colframe=accent!28,
  colbacktitle=accent!10,
  coltitle=accent!90!black,
  fonttitle=\bfseries,
  borderline west={0.9pt}{0pt}{accent!45},
  title={#1},
  before skip=6pt,
  after skip=8pt,
}

\newtcolorbox{ideabox}[1]{
  enhanced,
  colback=green!3,
  colframe=accent!20,
  colbacktitle=softgray!15,
  coltitle=accent!90!black,
  fonttitle=\bfseries,
  title={#1},
  boxed title style={
    sharp corners,
    boxrule=0pt,
    frame empty,
    interior style={fill=softgray!15},
    top=2pt,bottom=2pt,left=6pt,right=6pt,
  },
  borderline west={0.9pt}{0pt}{accent!35},
  before skip=6pt,
  after skip=8pt,
}

\newtcolorbox{warningbox}[1]{
  enhanced,
  colback=red!3,
  colframe=red!25,
  colbacktitle=red!10,
  coltitle=red!70!black,
  fonttitle=\bfseries,
  borderline west={0.9pt}{0pt}{red!35},
  title={#1},
  before skip=6pt,
  after skip=8pt,
}

\newtcolorbox{abstractbox}{
  enhanced,
  colback=softback,
  colframe=accent!28,
  coltitle=accent!90!black,
  fonttitle=\bfseries,
  title=摘要,
  boxrule=0.55pt,
  arc=2pt,
  left=8pt,right=8pt,top=6pt,bottom=6pt,
  before skip=4pt,
  after skip=10pt,
}

\DefineVerbatimEnvironment{codeblock}{Verbatim}{
  breaklines=true,
  breakanywhere=true,
  fontsize=\footnotesize,
  frame=single,
  rulecolor=\color{accent!28},
  framesep=3mm,
  baselinestretch=1,
}

\newcommand{\N}{\mathbb{N}}
\newcommand{\Z}{\mathbb{Z}}
\newcommand{\Q}{\mathbb{Q}}
\newcommand{\R}{\mathbb{R}}
\newcommand{\C}{\mathbb{C}}
\newcommand{\eps}{\varepsilon}
\newcommand{\id}{\mathrm{id}}

\title{\bfseries\color{accent}{LaTeX 讲义格式速记}\\[3pt]
\large\color{softgray}{复现本次版式}}
\author{}
\date{\color{softgray}{\today}}

\begin{document}
\maketitle

\begin{abstractbox}
这份文件把 \texttt{latex-style-notes.md} 直接改写成可编译的 LaTeX 版本，并且正文本身就使用同一套版式。以后新建讲义时，保留前导区与盒子定义，再替换正文即可。
\end{abstractbox}

\begin{multicols}{2}

\section{模板核心}

\begin{itemize}
\item 文档类：\texttt{ctexart}，参数为 \texttt{11pt,a4paper}。
\item 版面：\texttt{geometry[margin=1.7cm]}，两栏环境使用 \texttt{multicol}。
\item 列间距：\texttt{\textbackslash columnsep=0.8cm}。
\item 段落：\texttt{\textbackslash parindent=2em}，\texttt{\textbackslash parskip=0.3em}。
\item 常用包集中放在导言区，建议保持下面这组基础组合。
\end{itemize}

\begin{codeblock}
\usepackage{amsmath,amssymb,amsthm,mathtools,bm}
\usepackage{enumitem}
\usepackage{hyperref}
\usepackage{multicol}
\usepackage[most]{tcolorbox}
\usepackage{tikz}
\usepackage{microtype}
\usepackage{titlesec}
\usepackage{cleveref}
\end{codeblock}

\subsection{颜色定义}

\begin{codeblock}
\definecolor{accent}{RGB}{39,113,172}
\definecolor{softback}{RGB}{246,249,253}
\definecolor{softgray}{RGB}{110,120,135}
\end{codeblock}

\subsection{超链接与列表}

\begin{itemize}
\item 超链接建议启用 \texttt{colorlinks=true}，并统一用 \texttt{accent!80!black}。
\item \texttt{itemize} 左缩进设为 \texttt{1.2em}，\texttt{enumerate} 左缩进设为 \texttt{1.4em}。
\item 两类列表都把 \texttt{itemsep} 压到 \texttt{0.2em}，适合讲义型紧凑排版。
\end{itemize}

\begin{ideabox}{最小可复用骨架}
如果只想快速迁移风格，至少保留四组设置：页面几何、颜色定义、标题样式、盒子体系。正文内容可以完全替换。
\end{ideabox}

\section{标题样式}

\begin{codeblock}
\titleformat{\section}{\large\bfseries\color{accent}}{\thesection}{0.6em}{}[\titlerule]
\titleformat{\subsection}{\normalsize\bfseries\color{accent!90}}{\thesubsection}{0.5em}{}
\titleformat{\subsubsection}{\normalsize\bfseries\color{accent!80}}{\thesubsubsection}{0.5em}{}
\titlespacing*{\section}{0pt}{0.8em}{0.45em}
\titlespacing*{\subsection}{0pt}{0.6em}{0.3em}
\titlespacing*{\subsubsection}{0pt}{0.5em}{0.2em}
\end{codeblock}

\begin{examplebox}{标题示例}
主标题用 \texttt{accent}，副标题改用 \texttt{softgray}，这样层级清楚，而且不会比正文更重。讲义封面通常只保留主标题和副标题两层信息即可。
\end{examplebox}

\begin{codeblock}
\title{\bfseries\color{accent}{主标题} \\[3pt]
\large\color{softgray}{副标题}}
\end{codeblock}

\section{盒子体系（tcolorbox）}

\subsection{全局设置}

\begin{codeblock}
\tcbset{
  boxrule=0.55pt,
  arc=2pt,
  left=6pt,right=6pt,top=5pt,bottom=5pt,
  colframe=accent!25,
  colback=softback,
}
\end{codeblock}

\subsection{示例盒子}

\begin{examplebox}{examplebox}
适合放例题、小模板、定义后的直接应用，边线和底色都保持轻量，不抢正文。
\end{examplebox}

\begin{codeblock}
\newtcolorbox{examplebox}[1]{
  enhanced,
  colback=accent!3,
  colframe=accent!28,
  colbacktitle=accent!10,
  coltitle=accent!90!black,
  fonttitle=\bfseries,
  borderline west={0.9pt}{0pt}{accent!45},
  title={#1},
  before skip=6pt,
  after skip=8pt,
}
\end{codeblock}

\subsection{引导思考盒子}

\begin{ideabox}{ideabox}
适合放“怎么想”“怎么选方法”“做题先看什么”这类元信息。它的标题条更柔和，读者不会把它误认为正式定理。
\end{ideabox}

\begin{codeblock}
\newtcolorbox{ideabox}[1]{
  enhanced,
  colback=green!3,
  colframe=accent!20,
  colbacktitle=softgray!15,
  coltitle=accent!90!black,
  fonttitle=\bfseries,
  title={#1},
  boxed title style={
    sharp corners,
    boxrule=0pt,
    frame empty,
    interior style={fill=softgray!15},
    top=2pt,bottom=2pt,left=6pt,right=6pt,
  },
  borderline west={0.9pt}{0pt}{accent!35},
  before skip=6pt,
  after skip=8pt,
}
\end{codeblock}

\subsection{警告盒子}

\begin{warningbox}{warningbox}
适合放易错点、编译提醒、定义域条件或常见逻辑漏洞。颜色改成淡红后，读者扫页时能快速定位风险信息。
\end{warningbox}

\begin{codeblock}
\newtcolorbox{warningbox}[1]{
  enhanced,
  colback=red!3,
  colframe=red!25,
  colbacktitle=red!10,
  coltitle=red!70!black,
  fonttitle=\bfseries,
  borderline west={0.9pt}{0pt}{red!35},
  title={#1},
  before skip=6pt,
  after skip=8pt,
}
\end{codeblock}

\subsection{摘要盒子}

\begin{abstractbox}
摘要盒子适合放导读、文档说明、复习范围或使用方法。它的角色是开篇总览，不适合塞太多公式细节。
\end{abstractbox}

\begin{codeblock}
\newtcolorbox{abstractbox}{
  enhanced,
  colback=softback,
  colframe=accent!28,
  coltitle=accent!90!black,
  fonttitle=\bfseries,
  title=摘要,
  boxrule=0.55pt,
  arc=2pt,
  left=8pt,right=8pt,top=6pt,bottom=6pt,
  before skip=4pt,
  after skip=10pt,
}
\end{codeblock}

\section{宏与记号}

常用数学记号建议先统一定义，后续所有讲义都复用，能避免正文里反复手写格式。

\begin{codeblock}
\newcommand{\N}{\mathbb{N}}
\newcommand{\Z}{\mathbb{Z}}
\newcommand{\Q}{\mathbb{Q}}
\newcommand{\R}{\mathbb{R}}
\newcommand{\C}{\mathbb{C}}
\newcommand{\eps}{\varepsilon}
\newcommand{\id}{\mathrm{id}}
\end{codeblock}

\section{交叉引用与超链接}

\begin{itemize}
\item 引用章节、定理、公式时，建议统一交给 \texttt{cleveref}，例如 \texttt{\textbackslash Cref\{thm:prime-powers\}}。
\item 标题里一旦出现数学公式，就要配合 \texttt{texorpdfstring}，否则 PDF 书签容易报错或乱码。
\end{itemize}

\begin{codeblock}
\subsection{应用 1：快速求和
\texorpdfstring{$\sum_{n\le N}\varphi(n)$}{Sigma phi(n)}}
\end{codeblock}

\section{编译}

\begin{warningbox}{编译器选择}
推荐直接用 XeLaTeX。因为文档类是 \texttt{ctexart}，而正文里含中文，切到纯 pdfLaTeX 往往会直接卡在编码或字体上。
\end{warningbox}

\begin{codeblock}
xelatex -interaction=nonstopmode -halt-on-error yourfile.tex
\end{codeblock}

\section{复用方式}

\begin{enumerate}
\item 直接把“模板核心”“标题样式”“盒子体系”复制到新文档的导言区。
\item 把这份文件当成模板，在其基础上删除正文、替换标题和内容。
\item 如果想要单栏讲义，只需删掉 \texttt{multicol} 环境，其余颜色和盒子配置可以原样保留。
\end{enumerate}

\begin{ideabox}{建议的工作流}
先搭好导言区和盒子，再写摘要，再补各节内容。这样版式一次定死，后面只管填内容，不会在正文快写完时再返工排版。
\end{ideabox}

\end{multicols}

\end{document}
